\documentclass[11pt,a4paper]{article}

\usepackage{kotex}
\usepackage{amsmath,amssymb}
\usepackage{booktabs}
\usepackage{array}
\usepackage{geometry}
\usepackage{xcolor}
\usepackage{enumitem}
\usepackage{tcolorbox}
\usepackage{hyperref}
\usepackage{longtable}

\geometry{margin=25mm}

\definecolor{mainblue}{RGB}{45,72,96}
\definecolor{lightbluegray}{RGB}{232,238,243}
\definecolor{darkgray}{RGB}{55,60,65}
\definecolor{softgreen}{RGB}{230,240,235}
\definecolor{softyellow}{RGB}{247,242,225}

\hypersetup{
    colorlinks=true,
    linkcolor=mainblue,
    urlcolor=mainblue,
    citecolor=mainblue
}

\setlength{\parindent}{0pt}
\setlength{\parskip}{0.7em}

\setlist[itemize]{leftmargin=2em}
\setlist[enumerate]{leftmargin=2em}

\tcbset{
    colback=lightbluegray,
    colframe=mainblue,
    coltitle=white,
    colbacktitle=mainblue,
    fonttitle=\bfseries,
    arc=2mm,
    boxrule=0.7pt
}

\title{\textbf{Log-Weighted Cross-Entropy Framework for Imbalanced Classification}\\
\vspace{5mm}
\large 3차 논문 작성 피드백}
\author{고승호, 이현석}
\date{July 12, 2026}

\begin{document}

\maketitle

\section{전체 평가}

수정된 원고는 2차 피드백을 전반적으로 매우 잘 반영했음. 특히 Introduction에서 WCE의 Weight Explosion을 직접 수식으로 제시했고, LWCE, PLWCE, ES-LWCE의 역할을 서로 구분했으며, Related Work에는 Taxonomy와 Positioning Table이 들어가 논문의 위치가 명확해졌음. 또한 Section 3.2에서는 Monotonicity, Boundedness, Dynamic-range Compression, Gradient-scale Control을 Proposition 형태로 정리하여 이전보다 Theory Contribution이 크게 강화되었음.

현재 원고는 다음의 세 축이 비교적 선명해졌음.

\begin{center}
WCE의 Weight Explosion

$\Downarrow$

Logarithmic Compression

$\Downarrow$

Parameter-free LWCE와 두 Extension
\end{center}

특히 이전 피드백에서 강조했던 ``개별 Weight의 Boundedness''와 ``Rare-to-frequent Weight Ratio의 Logarithmic Compression''을 구분한 점이 좋음. ES-LWCE 역시 LWCE가 Positive Class Count에서 이미 bounded하다는 점을 인정하고, \(n_c=0\) 또는 Mini-batch-level Vanishing Count에 대한 Explicit Cap이라는 역할로 다시 설명한 점이 적절함.

\begin{tcolorbox}[title=종합 판단]
2차 피드백의 핵심은 대부분 잘 반영되었음. 현재까지는 Introduction, Related Work, Method Definition, Theoretical Properties의 흐름이 안정적임. 이제 논문의 완성도는 3.3 이후의 Method 세부화와 Section 4의 Experiment Design에 의해 결정될 것으로 판단됨.
\end{tcolorbox}

\section{이전 피드백 반영 여부}

\subsection{Boundedness 표현의 정밀화}

이전 피드백에서는 개별 Weight가 bounded하다는 사실과 Weight Ratio가 logarithmically compressed된다는 사실을 구분하라고 했음. 수정본은 이를 다음과 같이 정확히 구분함.

\[
0<w_c^{\mathrm{LWCE}}\leq \frac{1}{\log 2},
\qquad n_c\geq 1,
\]
반면
\[
\frac{w_{\min}^{\mathrm{LWCE}}}{w_{\max}^{\mathrm{LWCE}}}
=
\frac{\log(1+n_{\max})}{\log(1+n_{\min})}
=
O(\log \rho)
\]
로 설명하여, Ratio 자체는 constant-bounded가 아니라 WCE보다 훨씬 느리게 증가함을 명시함.

\begin{tcolorbox}[title=반영 평가]
매우 잘 반영됨. 이전 버전의 가장 큰 수학적 모호성이 해소되었음.
\end{tcolorbox}

\subsection{ES-LWCE의 역할 재정의}

수정본은 ES-LWCE의 역할을 \(n_c=0\)에서도 well-defined하고, Maximum Weight를 \(1/\epsilon\)으로 명시적으로 통제하며, Mini-batch-level Vanishing Count를 처리하는 Extension으로 정리함.

\begin{tcolorbox}[title=반영 평가]
잘 반영됨. ES-LWCE를 ``LWCE가 bounded하지 않아서 필요한 방법''이 아니라 ``Explicit, User-controlled Cap''을 제공하는 방법으로 Positioning한 점이 좋음.
\end{tcolorbox}

\subsection{PLWCE와 ES-LWCE의 Orthogonality}

수정본은 PLWCE가 Focusing Strength를 조절하고, ES-LWCE는 Maximum Weight와 Extreme-tail Smoothing을 담당한다고 명시함. 또한 두 방법이 결합 가능하지만 각 효과를 분리해서 보기 위해 독립적으로 평가한다고 설명함.

\begin{tcolorbox}[title=반영 평가]
잘 반영됨. 다만 실제 Experiment에서 이 설명을 뒷받침하는 Ablation이 반드시 필요함.
\end{tcolorbox}

\subsection{Theoretical Properties의 강화}

현재는 Monotonicity, Boundedness, Dynamic-range Compression, Gradient-scale Control의 네 개 Proposition으로 확장되었음. 특히 Proposition 4는 Weight Compression이 Optimization Stability로 이어지는 Mechanism을 설명한다는 점에서 중요함.

\begin{tcolorbox}[title=반영 평가]
매우 잘 반영됨. Theory Contribution이 독립적인 논문 기여로 보일 수 있을 정도로 개선되었음.
\end{tcolorbox}

\section{현재까지의 주요 보완점}

\subsection{Gradient-scale Control의 표현을 조금 더 엄밀하게 다듬을 것}

실제 Gradient Magnitude는 Prediction Probability와 Feature Representation에도 의존함. 따라서 Minority-to-majority Gradient Ratio가 Weight Ratio를 그대로 ``inherits''한다고 표현하면 다소 강하게 들릴 수 있음.

\begin{tcolorbox}[title=수정 권장]
``the class-dependent multiplicative component of the gradient scale follows the same asymptotic order as the corresponding weight ratio''처럼 표현하는 것이 더 안전함.
\end{tcolorbox}

\subsection{\(O(\log \rho)\) 표현의 조건을 명시할 것}

Dynamic-range Compression은 \(n_{\min}\)이 고정되거나 Lower-bounded된 상황을 암묵적으로 사용함.

\begin{tcolorbox}[title=수정 권장]
``For fixed \(n_{\min}\geq1\), the LWCE ratio grows as \(O(\log \rho)\)''와 같이 조건을 명시하는 것이 좋음.
\end{tcolorbox}

\subsection{Weight Normalization 여부를 명확히 해야 함}

실제 학습에서 Raw Weight를 사용하는지, Mean-normalized Weight를 사용하는지 반드시 명시해야 함. 예를 들어
\[
\widetilde{w}_c
=
\frac{Cw_c}{\sum_{j=1}^{C}w_j}
\]
와 같은 Normalization을 사용할 수 있음.

\begin{tcolorbox}[title=중요한 질문]
Baseline WCE와 CB Loss도 동일한 Scale Convention으로 비교해야 공정함.
\end{tcolorbox}

\section{3.3 이후 작성 방향}

\subsection{3.3 PLWCE Parameter Selection}

3.3절은 PLWCE의 \(\alpha\) 선택 절차를 재현 가능하게 설명해야 함.

\begin{itemize}
\item Candidate Set
\item Proxy Subset 구성 방식
\item Validation Metric
\item Selection Frequency
\item Additional Computational Cost
\item Test Data 미사용 명시
\end{itemize}

또한 \(\alpha\)의 의미를 다음과 같이 설명하면 좋음.

\begin{itemize}
\item \(0<\alpha<1\): Mild Re-weighting
\item \(\alpha=1\): LWCE
\item \(\alpha>1\): Strong Tail Emphasis
\end{itemize}

\begin{tcolorbox}[title=주의점]
Proxy Subset이 Training Data의 일부인지 별도의 Validation Set인지 명확히 해야 하며, Test Set을 이용한 \(\alpha\) 선택은 허용되지 않음.
\end{tcolorbox}

\subsection{3.4 ES-LWCE and Practical Use Cases}

ES-LWCE는 Mini-batch-level Class Frequency 또는 Dynamic/Streaming Class Count를 사용하는 경우에 초점을 맞추는 것이 좋음.

\[
w_c^{\mathrm{ES}}
=
\frac{1}{\log(1+n_c)+\epsilon}
\]
에서
\[
w_c^{\mathrm{ES}}\leq \frac{1}{\epsilon}
\]
이므로 \(\epsilon\)은 다음 Trade-off를 가짐.

\begin{itemize}
\item Small \(\epsilon\): Strong Minority Emphasis
\item Large \(\epsilon\): Conservative Gradient Scaling
\end{itemize}

\begin{tcolorbox}[title=개선 방향]
ES-LWCE를 Dataset-level Count에서 반드시 필요한 방법처럼 설명하기보다, Vanishing Count 또는 Batch-adaptive Count가 가능한 상황에서 유용한 Extension으로 Positioning하는 것이 좋음.
\end{tcolorbox}

\subsection{3.5 Computational Complexity and Implementation}

제안 방법의 장점 중 하나는 계산량이 거의 증가하지 않는다는 점임.

\begin{itemize}
\item Weight는 Class Count로부터 한 번 계산 가능함.
\item Per-sample Cost는 기존 WCE와 동일한 수준임.
\item Additional Model Parameter가 없음.
\item PLWCE의 Grid Search만 제한적인 추가 비용을 요구함.
\end{itemize}

추천 문장은 다음과 같음.

\begin{quote}
Once the class weights are precomputed, LWCE and ES-LWCE incur the same asymptotic training complexity as WCE, requiring only a scalar multiplication per sample.
\end{quote}

\subsection{그래프 추가}

parameter 가 변화하면서 loss function 이 어떻게 달라지는 지 등등 그래프를 적절히 추가하면 리뷰어들의 이해를 도울 수 있을 것. (학술대회 발표자료 참고)

\section{Section 4 Experiment 작성 방향}

\subsection{Experiment의 핵심 질문}

\begin{enumerate}
\item Imbalance가 심해질수록 LWCE의 장점이 커지는가?
\item LWCE는 WCE의 Optimization Instability를 실제로 줄이는가?
\item PLWCE의 \(\alpha\)는 Tail Emphasis를 의도한 대로 조절하는가?
\item ES-LWCE는 Vanishing Count 또는 Extreme Rare Class에서 실제로 안정적인가?
\end{enumerate}

\subsection{Baseline}

최소한 CE, WCE, Focal Loss, CB Loss, LDAM 또는 Logit Adjustment, LWCE, PLWCE, ES-LWCE를 포함하는 것이 좋음.

\subsection{평가지표}

Overall Accuracy만 사용하면 안 되며 다음 지표를 함께 보고해야 함.

\begin{itemize}
\item Macro-F1
\item Balanced Accuracy
\item Geometric Mean
\item Minority-class Recall
\item Worst-class Accuracy
\item Per-class Accuracy
\end{itemize}

\begin{tcolorbox}[title=핵심]
Imbalanced Classification 논문에서는 Macro-F1과 Worst-class Accuracy를 핵심 지표로 두는 것이 좋음.
\end{tcolorbox}

\subsection{Ablation Study}

\begin{center}
\begin{tabular}{lccc}
\toprule
Method & Log Compression & Power \(\alpha\) & Smoothing \(\epsilon\) \\
\midrule
WCE & X & X & X \\
LWCE & O & X & X \\
PLWCE & O & O & X \\
ES-LWCE & O & X & O \\
Combined & O & O & O \\
\bottomrule
\end{tabular}
\end{center}

Combined Variant를 Main Method로 제안하지 않더라도 Supplementary Ablation으로 넣으면 Orthogonality 주장을 검증할 수 있음.

\subsection{Sensitivity Analysis}

\(\alpha\)와 \(\epsilon\) 변화에 따른 Macro-F1, Worst-class Accuracy, Gradient Norm을 그래프로 보여주는 것이 좋음.

\subsection{Optimization Stability 분석}

논문의 핵심 Claim이 Weight Explosion 완화이므로 다음 시각화가 중요함.

\begin{itemize}
\item Epoch별 Training Loss
\item Epoch별 Gradient Norm
\item Minority/Majority Gradient Ratio
\item Class Weight Distribution
\end{itemize}

\begin{tcolorbox}[title=중요]
단순 성능 향상만으로는 ``Weight Explosion을 완화했다''는 Mechanism을 충분히 증명하기 어려움. Gradient Norm과 Weight Distribution 분석을 반드시 포함하는 것이 좋음.
\end{tcolorbox}

\subsection{Statistical Analysis}

Dataset이 여러 개이므로 경우 다음을 추천함 (태욱, 인준방법 말고 기존 방법).

\begin{itemize}
\item Friedman Test
\item Post-hoc Holm Procedure
\item Mean Rank
\item Win/Tie/Loss Count
\end{itemize}

\section{수정 우선순위}

\subsection{Priority 1}

\begin{enumerate}
\item Weight Normalization 여부 명시
\item Proposition 4 표현 정밀화
\item 3.3에서 PLWCE Selection Procedure 구체화
\item 3.4에서 ES-LWCE의 Practical Use Case 명확화
\item Experiment에 Macro-F1, Worst-class Accuracy, Gradient Norm 포함
\end{enumerate}

\subsection{Priority 2}

\begin{enumerate}
\item \(\alpha\), \(\epsilon\) Sensitivity Analysis
\item Combined Variant Ablation
\item Computational Cost 비교
\item Statistical Significance Test
\end{enumerate}


\end{document}
